% Two-tank adaptation of the existing three-tank worksheet figure.
% Load once in the preamble: % Two-tank adaptation of the existing three-tank worksheet figure.
% Load once in the preamble: % Two-tank adaptation of the existing three-tank worksheet figure.
% Load once in the preamble: % Two-tank adaptation of the existing three-tank worksheet figure.
% Load once in the preamble: \input{TwoTankDiagram.tex}
% Then use \TwoTankDiagram{key=value,...}. All keys are optional.
% capacity1/2 and volume1/2: L. initialSalt1/2: g.
% inflow, flow12, flow21, outflow: L/min. inflowconcentration: g/L.
% saltRate1/2: dry salt in g/min; zero hides the conveyor.
% Liquid height is proportional to volume/capacity. Compact headings retain
% the original T_i (capacity L) format, without a verbose capacity label.
% Requires 0 < volume <= capacity and nonnegative flow/source data.
% Drawing style and conveyor geometry are inherited from TankDiagram.tex.
\usepackage{tikz}
\usetikzlibrary{arrows.meta}
\providecolor{mtuGreen}{RGB}{0,79,42}
\providecolor{mtuGreenLight}{RGB}{60,133,90}
\providecolor{mtuGreenPale}{RGB}{223,236,227}
\pgfkeys{/two tank worksheet/.is family,/two tank worksheet,
 capacity1/.store in=\TTcapone,capacity2/.store in=\TTcaptwo,
 volume1/.store in=\TTvolone,volume2/.store in=\TTvoltwo,
 initialSalt1/.store in=\TTsaltone,initialSalt2/.store in=\TTsalttwo,
 inflow/.store in=\TTqin,inflowconcentration/.store in=\TTcin,
 flow12/.store in=\TTqonetwo,flow21/.store in=\TTqtwoone,
 outflow/.store in=\TTqout,
 saltRate1/.store in=\TTdryone,saltRate2/.store in=\TTdrytwo}
% Shared compact labels used by both lectures and worksheets.
\newcommand{\TwoTankUnitHeading}[2]{%
  \node[font=\bfseries] at (0,3.12) {$T_{#1}$ (#2 L)};}
\newcommand{\TwoTankUnitLabels}[4]{%
  \node[anchor=north,inner sep=1pt,font=\fontsize{7}{8}\selectfont] at (0,-.20)
  {$v_{#1}(0)=#3\,\mathrm{L},\;s_{#1}(0)= #4\,\mathrm{g}$};}
% One cylindrical tank: index, horizontal position, capacity, initial volume,
% initial salt amount in g. The liquid height follows initial volume/capacity.
\newcommand{\TwoTankUnit}[5]{%
  \begin{scope}[shift={(#2,0)}]
    \pgfmathsetmacro{\TwoTankLevel}{2.6*(#4)/(#3)}
    \path[fill=mtuGreenPale] (-.9,0) arc[start angle=180,end angle=360,x radius=.9,y radius=.17]
      -- (.9,\TwoTankLevel) arc[start angle=0,end angle=180,x radius=.9,y radius=.17] -- cycle;
    \draw[mtuGreenLight,fill=mtuGreenLight!20] (0,\TwoTankLevel) ellipse[x radius=.9,y radius=.17];
    \draw[thick] (-.9,2.6)--(-.9,0) arc[start angle=180,end angle=360,x radius=.9,y radius=.17] -- (.9,2.6);
    \draw[thick] (0,2.6) ellipse[x radius=.9,y radius=.17];
    \TwoTankUnitHeading{#1}{#3}
    \TwoTankUnitLabels{#1}{#3}{#4}{#5}
    \coordinate (tank#1wh) at (-.9,1.95);
    \coordinate (tank#1eh) at (.9,1.95);
    \coordinate (tank#1wl) at (-.9,.7);
    \coordinate (tank#1el) at (.9,.7);
  \end{scope}}

% Optional independent continuous salt rates in g/min. Zero disables each source.
% Add these rates to the salt-input vector, not to initial salt masses.
\newcommand{\TwoTankUnitSaltSource}[2]{%
  \ifdim#2pt>0pt\relax
    \begin{scope}[shift={(#1,0)},xshift=0.15cm,yshift=-0.12cm]
      % Belt sits below the capacity heading and delivers inward over the rim.
      \node[anchor=west,font=\small] at (1.12,3.20) {$#2\,\mathrm{g/min}$};
      \draw[thick,mtuGreen,rounded corners=2pt,fill=white]
        (0.48,2.80) rectangle (1.92,2.96);
      \draw[mtuGreen] (0.60,2.88) circle[radius=0.055];
      \draw[mtuGreen] (1.80,2.88) circle[radius=0.055];
      \draw[-{Stealth[length=1.5mm]},mtuGreen] (1.55,2.88)--(0.90,2.88);
      \foreach \sx/\sy in {1.65/3.00,1.43/3.00,1.20/3.00,0.98/3.00,0.75/3.00,0.50/2.75,0.48/2.64,0.53/2.53,0.48/2.42}{%
        \fill[red] (\sx,\sy) circle[radius=0.022];}
    \end{scope}
  \fi}

\newcommand{\TwoTankDiagram}[1]{%
 \begingroup
 \pgfkeys{/two tank worksheet/.cd,
 capacity1=100,capacity2=120,volume1=60,volume2=80,
 initialSalt1=120,initialSalt2=80,
 inflow=3,inflowconcentration=3,flow12=5,flow21=2,outflow=3,
 saltRate1=0,saltRate2=0,#1}
 \begin{tikzpicture}[font=\small,
 pipe/.style={-{Stealth[length=2.5mm,width=1.7mm]},thick,mtuGreen},
 rate/.style={above,align=center,text=black,fill=white,inner sep=3pt}]
 \TwoTankUnit{1}{0}{\TTcapone}{\TTvolone}{\TTsaltone}
 \TwoTankUnit{2}{4.6}{\TTcaptwo}{\TTvoltwo}{\TTsalttwo}
 \TwoTankUnitSaltSource{0}{\TTdryone}
 \TwoTankUnitSaltSource{4.6}{\TTdrytwo}
 \draw[pipe] (-3.7,1.95)--(tank1wh) node[midway,rate] {$\begin{aligned}q_{\mathrm{in}}&=\TTqin\ \mathrm{L/min}\\c_{\mathrm{in}}&=\TTcin\ \mathrm{g/L}\end{aligned}$};
 \draw[pipe] (tank1eh)--(tank2wh) node[midway,rate] {$q_{12}=\TTqonetwo\ \mathrm{L/min}$};
 \draw[pipe] (tank2wl)--(tank1el) node[midway,rate] {$q_{21}=\TTqtwoone\ \mathrm{L/min}$};
 \draw[pipe] (tank2el)--(8.3,.7) node[midway,rate] {$q_{\mathrm{out}}=\TTqout\ \mathrm{L/min}$};
 \end{tikzpicture}%
 \endgroup}

% Then use \TwoTankDiagram{key=value,...}. All keys are optional.
% capacity1/2 and volume1/2: L. initialSalt1/2: g.
% inflow, flow12, flow21, outflow: L/min. inflowconcentration: g/L.
% saltRate1/2: dry salt in g/min; zero hides the conveyor.
% Liquid height is proportional to volume/capacity. Compact headings retain
% the original T_i (capacity L) format, without a verbose capacity label.
% Requires 0 < volume <= capacity and nonnegative flow/source data.
% Drawing style and conveyor geometry are inherited from TankDiagram.tex.
\usepackage{tikz}
\usetikzlibrary{arrows.meta}
\providecolor{mtuGreen}{RGB}{0,79,42}
\providecolor{mtuGreenLight}{RGB}{60,133,90}
\providecolor{mtuGreenPale}{RGB}{223,236,227}
\pgfkeys{/two tank worksheet/.is family,/two tank worksheet,
 capacity1/.store in=\TTcapone,capacity2/.store in=\TTcaptwo,
 volume1/.store in=\TTvolone,volume2/.store in=\TTvoltwo,
 initialSalt1/.store in=\TTsaltone,initialSalt2/.store in=\TTsalttwo,
 inflow/.store in=\TTqin,inflowconcentration/.store in=\TTcin,
 flow12/.store in=\TTqonetwo,flow21/.store in=\TTqtwoone,
 outflow/.store in=\TTqout,
 saltRate1/.store in=\TTdryone,saltRate2/.store in=\TTdrytwo}
% Shared compact labels used by both lectures and worksheets.
\newcommand{\TwoTankUnitHeading}[2]{%
  \node[font=\bfseries] at (0,3.12) {$T_{#1}$ (#2 L)};}
\newcommand{\TwoTankUnitLabels}[4]{%
  \node[anchor=north,inner sep=1pt,font=\fontsize{7}{8}\selectfont] at (0,-.20)
  {$v_{#1}(0)=#3\,\mathrm{L},\;s_{#1}(0)= #4\,\mathrm{g}$};}
% One cylindrical tank: index, horizontal position, capacity, initial volume,
% initial salt amount in g. The liquid height follows initial volume/capacity.
\newcommand{\TwoTankUnit}[5]{%
  \begin{scope}[shift={(#2,0)}]
    \pgfmathsetmacro{\TwoTankLevel}{2.6*(#4)/(#3)}
    \path[fill=mtuGreenPale] (-.9,0) arc[start angle=180,end angle=360,x radius=.9,y radius=.17]
      -- (.9,\TwoTankLevel) arc[start angle=0,end angle=180,x radius=.9,y radius=.17] -- cycle;
    \draw[mtuGreenLight,fill=mtuGreenLight!20] (0,\TwoTankLevel) ellipse[x radius=.9,y radius=.17];
    \draw[thick] (-.9,2.6)--(-.9,0) arc[start angle=180,end angle=360,x radius=.9,y radius=.17] -- (.9,2.6);
    \draw[thick] (0,2.6) ellipse[x radius=.9,y radius=.17];
    \TwoTankUnitHeading{#1}{#3}
    \TwoTankUnitLabels{#1}{#3}{#4}{#5}
    \coordinate (tank#1wh) at (-.9,1.95);
    \coordinate (tank#1eh) at (.9,1.95);
    \coordinate (tank#1wl) at (-.9,.7);
    \coordinate (tank#1el) at (.9,.7);
  \end{scope}}

% Optional independent continuous salt rates in g/min. Zero disables each source.
% Add these rates to the salt-input vector, not to initial salt masses.
\newcommand{\TwoTankUnitSaltSource}[2]{%
  \ifdim#2pt>0pt\relax
    \begin{scope}[shift={(#1,0)},xshift=0.15cm,yshift=-0.12cm]
      % Belt sits below the capacity heading and delivers inward over the rim.
      \node[anchor=west,font=\small] at (1.12,3.20) {$#2\,\mathrm{g/min}$};
      \draw[thick,mtuGreen,rounded corners=2pt,fill=white]
        (0.48,2.80) rectangle (1.92,2.96);
      \draw[mtuGreen] (0.60,2.88) circle[radius=0.055];
      \draw[mtuGreen] (1.80,2.88) circle[radius=0.055];
      \draw[-{Stealth[length=1.5mm]},mtuGreen] (1.55,2.88)--(0.90,2.88);
      \foreach \sx/\sy in {1.65/3.00,1.43/3.00,1.20/3.00,0.98/3.00,0.75/3.00,0.50/2.75,0.48/2.64,0.53/2.53,0.48/2.42}{%
        \fill[red] (\sx,\sy) circle[radius=0.022];}
    \end{scope}
  \fi}

\newcommand{\TwoTankDiagram}[1]{%
 \begingroup
 \pgfkeys{/two tank worksheet/.cd,
 capacity1=100,capacity2=120,volume1=60,volume2=80,
 initialSalt1=120,initialSalt2=80,
 inflow=3,inflowconcentration=3,flow12=5,flow21=2,outflow=3,
 saltRate1=0,saltRate2=0,#1}
 \begin{tikzpicture}[font=\small,
 pipe/.style={-{Stealth[length=2.5mm,width=1.7mm]},thick,mtuGreen},
 rate/.style={above,align=center,text=black,fill=white,inner sep=3pt}]
 \TwoTankUnit{1}{0}{\TTcapone}{\TTvolone}{\TTsaltone}
 \TwoTankUnit{2}{4.6}{\TTcaptwo}{\TTvoltwo}{\TTsalttwo}
 \TwoTankUnitSaltSource{0}{\TTdryone}
 \TwoTankUnitSaltSource{4.6}{\TTdrytwo}
 \draw[pipe] (-3.7,1.95)--(tank1wh) node[midway,rate] {$\begin{aligned}q_{\mathrm{in}}&=\TTqin\ \mathrm{L/min}\\c_{\mathrm{in}}&=\TTcin\ \mathrm{g/L}\end{aligned}$};
 \draw[pipe] (tank1eh)--(tank2wh) node[midway,rate] {$q_{12}=\TTqonetwo\ \mathrm{L/min}$};
 \draw[pipe] (tank2wl)--(tank1el) node[midway,rate] {$q_{21}=\TTqtwoone\ \mathrm{L/min}$};
 \draw[pipe] (tank2el)--(8.3,.7) node[midway,rate] {$q_{\mathrm{out}}=\TTqout\ \mathrm{L/min}$};
 \end{tikzpicture}%
 \endgroup}

% Then use \TwoTankDiagram{key=value,...}. All keys are optional.
% capacity1/2 and volume1/2: L. initialSalt1/2: g.
% inflow, flow12, flow21, outflow: L/min. inflowconcentration: g/L.
% saltRate1/2: dry salt in g/min; zero hides the conveyor.
% Liquid height is proportional to volume/capacity. Compact headings retain
% the original T_i (capacity L) format, without a verbose capacity label.
% Requires 0 < volume <= capacity and nonnegative flow/source data.
% Drawing style and conveyor geometry are inherited from TankDiagram.tex.
\usepackage{tikz}
\usetikzlibrary{arrows.meta}
\providecolor{mtuGreen}{RGB}{0,79,42}
\providecolor{mtuGreenLight}{RGB}{60,133,90}
\providecolor{mtuGreenPale}{RGB}{223,236,227}
\pgfkeys{/two tank worksheet/.is family,/two tank worksheet,
 capacity1/.store in=\TTcapone,capacity2/.store in=\TTcaptwo,
 volume1/.store in=\TTvolone,volume2/.store in=\TTvoltwo,
 initialSalt1/.store in=\TTsaltone,initialSalt2/.store in=\TTsalttwo,
 inflow/.store in=\TTqin,inflowconcentration/.store in=\TTcin,
 flow12/.store in=\TTqonetwo,flow21/.store in=\TTqtwoone,
 outflow/.store in=\TTqout,
 saltRate1/.store in=\TTdryone,saltRate2/.store in=\TTdrytwo}
% Shared compact labels used by both lectures and worksheets.
\newcommand{\TwoTankUnitHeading}[2]{%
  \node[font=\bfseries] at (0,3.12) {$T_{#1}$ (#2 L)};}
\newcommand{\TwoTankUnitLabels}[4]{%
  \node[anchor=north,inner sep=1pt,font=\fontsize{7}{8}\selectfont] at (0,-.20)
  {$v_{#1}(0)=#3\,\mathrm{L},\;s_{#1}(0)= #4\,\mathrm{g}$};}
% One cylindrical tank: index, horizontal position, capacity, initial volume,
% initial salt amount in g. The liquid height follows initial volume/capacity.
\newcommand{\TwoTankUnit}[5]{%
  \begin{scope}[shift={(#2,0)}]
    \pgfmathsetmacro{\TwoTankLevel}{2.6*(#4)/(#3)}
    \path[fill=mtuGreenPale] (-.9,0) arc[start angle=180,end angle=360,x radius=.9,y radius=.17]
      -- (.9,\TwoTankLevel) arc[start angle=0,end angle=180,x radius=.9,y radius=.17] -- cycle;
    \draw[mtuGreenLight,fill=mtuGreenLight!20] (0,\TwoTankLevel) ellipse[x radius=.9,y radius=.17];
    \draw[thick] (-.9,2.6)--(-.9,0) arc[start angle=180,end angle=360,x radius=.9,y radius=.17] -- (.9,2.6);
    \draw[thick] (0,2.6) ellipse[x radius=.9,y radius=.17];
    \TwoTankUnitHeading{#1}{#3}
    \TwoTankUnitLabels{#1}{#3}{#4}{#5}
    \coordinate (tank#1wh) at (-.9,1.95);
    \coordinate (tank#1eh) at (.9,1.95);
    \coordinate (tank#1wl) at (-.9,.7);
    \coordinate (tank#1el) at (.9,.7);
  \end{scope}}

% Optional independent continuous salt rates in g/min. Zero disables each source.
% Add these rates to the salt-input vector, not to initial salt masses.
\newcommand{\TwoTankUnitSaltSource}[2]{%
  \ifdim#2pt>0pt\relax
    \begin{scope}[shift={(#1,0)},xshift=0.15cm,yshift=-0.12cm]
      % Belt sits below the capacity heading and delivers inward over the rim.
      \node[anchor=west,font=\small] at (1.12,3.20) {$#2\,\mathrm{g/min}$};
      \draw[thick,mtuGreen,rounded corners=2pt,fill=white]
        (0.48,2.80) rectangle (1.92,2.96);
      \draw[mtuGreen] (0.60,2.88) circle[radius=0.055];
      \draw[mtuGreen] (1.80,2.88) circle[radius=0.055];
      \draw[-{Stealth[length=1.5mm]},mtuGreen] (1.55,2.88)--(0.90,2.88);
      \foreach \sx/\sy in {1.65/3.00,1.43/3.00,1.20/3.00,0.98/3.00,0.75/3.00,0.50/2.75,0.48/2.64,0.53/2.53,0.48/2.42}{%
        \fill[red] (\sx,\sy) circle[radius=0.022];}
    \end{scope}
  \fi}

\newcommand{\TwoTankDiagram}[1]{%
 \begingroup
 \pgfkeys{/two tank worksheet/.cd,
 capacity1=100,capacity2=120,volume1=60,volume2=80,
 initialSalt1=120,initialSalt2=80,
 inflow=3,inflowconcentration=3,flow12=5,flow21=2,outflow=3,
 saltRate1=0,saltRate2=0,#1}
 \begin{tikzpicture}[font=\small,
 pipe/.style={-{Stealth[length=2.5mm,width=1.7mm]},thick,mtuGreen},
 rate/.style={above,align=center,text=black,fill=white,inner sep=3pt}]
 \TwoTankUnit{1}{0}{\TTcapone}{\TTvolone}{\TTsaltone}
 \TwoTankUnit{2}{4.6}{\TTcaptwo}{\TTvoltwo}{\TTsalttwo}
 \TwoTankUnitSaltSource{0}{\TTdryone}
 \TwoTankUnitSaltSource{4.6}{\TTdrytwo}
 \draw[pipe] (-3.7,1.95)--(tank1wh) node[midway,rate] {$\begin{aligned}q_{\mathrm{in}}&=\TTqin\ \mathrm{L/min}\\c_{\mathrm{in}}&=\TTcin\ \mathrm{g/L}\end{aligned}$};
 \draw[pipe] (tank1eh)--(tank2wh) node[midway,rate] {$q_{12}=\TTqonetwo\ \mathrm{L/min}$};
 \draw[pipe] (tank2wl)--(tank1el) node[midway,rate] {$q_{21}=\TTqtwoone\ \mathrm{L/min}$};
 \draw[pipe] (tank2el)--(8.3,.7) node[midway,rate] {$q_{\mathrm{out}}=\TTqout\ \mathrm{L/min}$};
 \end{tikzpicture}%
 \endgroup}

% Then use \TwoTankDiagram{key=value,...}. All keys are optional.
% capacity1/2 and volume1/2: L. initialSalt1/2: g.
% inflow, flow12, flow21, outflow: L/min. inflowconcentration: g/L.
% saltRate1/2: dry salt in g/min; zero hides the conveyor.
% Liquid height is proportional to volume/capacity. Compact headings retain
% the original T_i (capacity L) format, without a verbose capacity label.
% Requires 0 < volume <= capacity and nonnegative flow/source data.
% Drawing style and conveyor geometry are inherited from TankDiagram.tex.
\usepackage{tikz}
\usetikzlibrary{arrows.meta}
\providecolor{mtuGreen}{RGB}{0,79,42}
\providecolor{mtuGreenLight}{RGB}{60,133,90}
\providecolor{mtuGreenPale}{RGB}{223,236,227}
\pgfkeys{/two tank worksheet/.is family,/two tank worksheet,
 capacity1/.store in=\TTcapone,capacity2/.store in=\TTcaptwo,
 volume1/.store in=\TTvolone,volume2/.store in=\TTvoltwo,
 initialSalt1/.store in=\TTsaltone,initialSalt2/.store in=\TTsalttwo,
 inflow/.store in=\TTqin,inflowconcentration/.store in=\TTcin,
 flow12/.store in=\TTqonetwo,flow21/.store in=\TTqtwoone,
 outflow/.store in=\TTqout,
 saltRate1/.store in=\TTdryone,saltRate2/.store in=\TTdrytwo}
% Shared compact labels used by both lectures and worksheets.
\newcommand{\TwoTankUnitHeading}[2]{%
  \node[font=\bfseries] at (0,3.12) {$T_{#1}$ (#2 L)};}
\newcommand{\TwoTankUnitLabels}[4]{%
  \node[anchor=north,inner sep=1pt,font=\fontsize{7}{8}\selectfont] at (0,-.20)
  {$v_{#1}(0)=#3\,\mathrm{L},\;s_{#1}(0)= #4\,\mathrm{g}$};}
% One cylindrical tank: index, horizontal position, capacity, initial volume,
% initial salt amount in g. The liquid height follows initial volume/capacity.
\newcommand{\TwoTankUnit}[5]{%
  \begin{scope}[shift={(#2,0)}]
    \pgfmathsetmacro{\TwoTankLevel}{2.6*(#4)/(#3)}
    \path[fill=mtuGreenPale] (-.9,0) arc[start angle=180,end angle=360,x radius=.9,y radius=.17]
      -- (.9,\TwoTankLevel) arc[start angle=0,end angle=180,x radius=.9,y radius=.17] -- cycle;
    \draw[mtuGreenLight,fill=mtuGreenLight!20] (0,\TwoTankLevel) ellipse[x radius=.9,y radius=.17];
    \draw[thick] (-.9,2.6)--(-.9,0) arc[start angle=180,end angle=360,x radius=.9,y radius=.17] -- (.9,2.6);
    \draw[thick] (0,2.6) ellipse[x radius=.9,y radius=.17];
    \TwoTankUnitHeading{#1}{#3}
    \TwoTankUnitLabels{#1}{#3}{#4}{#5}
    \coordinate (tank#1wh) at (-.9,1.95);
    \coordinate (tank#1eh) at (.9,1.95);
    \coordinate (tank#1wl) at (-.9,.7);
    \coordinate (tank#1el) at (.9,.7);
  \end{scope}}

% Optional independent continuous salt rates in g/min. Zero disables each source.
% Add these rates to the salt-input vector, not to initial salt masses.
\newcommand{\TwoTankUnitSaltSource}[2]{%
  \ifdim#2pt>0pt\relax
    \begin{scope}[shift={(#1,0)},xshift=0.15cm,yshift=-0.12cm]
      % Belt sits below the capacity heading and delivers inward over the rim.
      \node[anchor=west,font=\small] at (1.12,3.20) {$#2\,\mathrm{g/min}$};
      \draw[thick,mtuGreen,rounded corners=2pt,fill=white]
        (0.48,2.80) rectangle (1.92,2.96);
      \draw[mtuGreen] (0.60,2.88) circle[radius=0.055];
      \draw[mtuGreen] (1.80,2.88) circle[radius=0.055];
      \draw[-{Stealth[length=1.5mm]},mtuGreen] (1.55,2.88)--(0.90,2.88);
      \foreach \sx/\sy in {1.65/3.00,1.43/3.00,1.20/3.00,0.98/3.00,0.75/3.00,0.50/2.75,0.48/2.64,0.53/2.53,0.48/2.42}{%
        \fill[red] (\sx,\sy) circle[radius=0.022];}
    \end{scope}
  \fi}

\newcommand{\TwoTankDiagram}[1]{%
 \begingroup
 \pgfkeys{/two tank worksheet/.cd,
 capacity1=100,capacity2=120,volume1=60,volume2=80,
 initialSalt1=120,initialSalt2=80,
 inflow=3,inflowconcentration=3,flow12=5,flow21=2,outflow=3,
 saltRate1=0,saltRate2=0,#1}
 \begin{tikzpicture}[font=\small,
 pipe/.style={-{Stealth[length=2.5mm,width=1.7mm]},thick,mtuGreen},
 rate/.style={above,align=center,text=black,fill=white,inner sep=3pt}]
 \TwoTankUnit{1}{0}{\TTcapone}{\TTvolone}{\TTsaltone}
 \TwoTankUnit{2}{4.6}{\TTcaptwo}{\TTvoltwo}{\TTsalttwo}
 \TwoTankUnitSaltSource{0}{\TTdryone}
 \TwoTankUnitSaltSource{4.6}{\TTdrytwo}
 \draw[pipe] (-3.7,1.95)--(tank1wh) node[midway,rate] {$\begin{aligned}q_{\mathrm{in}}&=\TTqin\ \mathrm{L/min}\\c_{\mathrm{in}}&=\TTcin\ \mathrm{g/L}\end{aligned}$};
 \draw[pipe] (tank1eh)--(tank2wh) node[midway,rate] {$q_{12}=\TTqonetwo\ \mathrm{L/min}$};
 \draw[pipe] (tank2wl)--(tank1el) node[midway,rate] {$q_{21}=\TTqtwoone\ \mathrm{L/min}$};
 \draw[pipe] (tank2el)--(8.3,.7) node[midway,rate] {$q_{\mathrm{out}}=\TTqout\ \mathrm{L/min}$};
 \end{tikzpicture}%
 \endgroup}
